\documentclass[aps,prc,onecolumn,superscriptaddress,nofootinbib,10pt]{revtex4-2}
\usepackage{amsmath,amssymb}
\usepackage{booktabs}
\usepackage[colorlinks=true,linkcolor=blue,citecolor=blue,urlcolor=blue]{hyperref}

\newcommand{\nB}{n_B}
\newcommand{\nsat}{n_{\mathrm{sat}}}
\newcommand{\muB}{\mu_B}
\newcommand{\muC}{\mu_C}
\newcommand{\muS}{\mu_S}
\newcommand{\munue}{\mu_{\nu_e}}
\newcommand{\mueff}[1]{\mu^{\mathrm{eff}}_{#1}}
\newcommand{\meff}[1]{m^{*}_{#1}}

\begin{document}

\title{SFHo: a nonlinear relativistic mean-field model \\
       for hadronic matter --- as implemented in \texttt{eos/sfho}}

\author{M.~Guerrini}
\affiliation{University of Ferrara}

\begin{abstract}
The \texttt{eos/sfho} subpackage implements the SFHo parametrization of
Steiner, Hempel \& Fischer~\cite{Steiner2013}, a nonlinear relativistic
mean-field model in the Boguta--Bodmer family~\cite{BogutaBodmer1977} with the
isoscalar--isovector cross coupling of Steiner, Prakash, Lattimer \&
Ellis~\cite{Steiner2005}, extended to the hyperon octet with the SFHoY and
SFHoY$^{*}$ couplings of Fortin, Oertel \& Provid\^encia~\cite{Fortin2017}, to
the $\Delta(1232)$ quartet, and to a thermal pseudoscalar meson gas following
Lavagno~\cite{Lavagno2010}. This note states the Lagrangian, the parameter
set, the field equations, the thermodynamics, and the closure of each
equilibrium mode the code supports. It is written from the code, as its
specification: every equation here is asserted by the \texttt{verify/} suite
or the test gates.

SFHo and DD2 are the two hadronic models of this repository and are
deliberately built the same way --- same module layout, same names, same
modes. The physics difference is one thing: DD2 puts the density dependence in
the \emph{couplings}, SFHo puts it in nonlinear \emph{self-interactions} of
the meson fields. SFHo therefore has constant $g_i$, no rearrangement
self-energy, and no \texttt{couplings.py}.
\end{abstract}

\maketitle

\section{Lagrangian and couplings}
\label{sec:lagrangian}

Baryons $i$ (nucleons; optionally the hyperon octet and the $\Delta$ quartet)
couple to the isoscalar--scalar $\sigma$, isoscalar--vector $\omega$,
isovector--vector $\rho$, and (for strange baryons) the hidden-strange vector
$\phi$, with \emph{constant} couplings $g_{Mi}$:
%
\begin{align}
\mathcal{L} ={}& \sum_i \bar\psi_i \Big[ \gamma_\mu \big( i\partial^\mu
      - g_{\omega i}\,\omega^\mu
      - g_{\rho i}\, I_{3i}\, \rho^\mu
      - g_{\phi i}\, \phi^\mu \big)
      - \big( m_i - g_{\sigma i}\,\sigma \big) \Big] \psi_i
      \nonumber\\
 &+ \tfrac12 \big( \partial_\mu \sigma\, \partial^\mu \sigma
      - m_\sigma^2 \sigma^2 \big) - U(\sigma)
   - \tfrac14 W_{\mu\nu} W^{\mu\nu} + \tfrac12 m_\omega^2\, \omega_\mu \omega^\mu
   + \tfrac{c_3}{4} (\omega_\mu \omega^\mu)^2
      \nonumber\\
 &- \tfrac14 R_{\mu\nu} R^{\mu\nu} + \tfrac12 m_\rho^2\, \rho_\mu \rho^\mu
   + \tfrac{c_4}{4} (\rho_\mu \rho^\mu)^2
   + A(\sigma, \omega)\, \rho_\mu \rho^\mu
   - \tfrac14 \Phi_{\mu\nu} \Phi^{\mu\nu} + \tfrac12 m_\phi^2\, \phi_\mu \phi^\mu ,
\label{eq:lagrangian}
\end{align}
%
where $I_{3i}$ is the third isospin component in the $\tau_3/2$ convention
($I_{3p} = +\tfrac12$, $I_{3n} = -\tfrac12$, $I_{3\Sigma^+} = +1$,
$I_{3\Delta^{++}} = +\tfrac32$), which is the convention the published SFHo
$\rho$ coupling is tabulated in. (DD2 uses $\tau_3 = \pm 1$ instead; the two
conventions absorb a factor of two into $g_\rho$, and neither model may be
read with the other's table.) Leptons ($e$) and, in trapped modes, electron
neutrinos enter as free Dirac gases; photons as a free Bose gas.

Three nonlinear terms carry what DD2 carries in $\Gamma_i(\nB)$. The scalar
self-interaction is Boguta--Bodmer~\cite{BogutaBodmer1977},
%
\begin{equation}
U(\sigma) = \frac{g_2}{3}\, \sigma^3 + \frac{g_3}{4}\, \sigma^4 ,
\qquad
\frac{\partial U}{\partial \sigma} = g_2 \sigma^2 + g_3 \sigma^3 ,
\label{eq:Usigma}
\end{equation}
%
and controls $K_{\mathrm{sat}}$ and $m^*/m$. The quartic vector terms $c_3,
c_4$ stiffen (or soften) the high-density vector sector. The third is the
isoscalar--isovector cross coupling of Steiner et al.~\cite{Steiner2005},
%
\begin{equation}
A(\sigma, \omega) = g_{\rho N}^2\, f(\sigma, \omega),
\qquad
f(\sigma, \omega) = \sum_{i=1}^{6} a_i\, \sigma^i
                  + \sum_{j=1}^{3} b_j\, \omega^{2j} ,
\label{eq:Afunction}
\end{equation}
%
which is what makes $L_{\mathrm{sym}}$ adjustable at fixed $E_{\mathrm{sym}}$
in this family. It generalises the $\Lambda_v \omega^2 \rho^2$ coupling of
Horowitz \& Piekarewicz~\cite{HorowitzPiekarewicz2001}, which is the single
term $b_1$: $A$ is a polynomial in $\sigma$ and in $\omega^2$ separately, so
$\partial^2 A/\partial\sigma\,\partial\omega \equiv 0$ --- a fact the analytic
Jacobian uses. The published values~\cite{Steiner2013}, as transcribed from
the CompOSE table, are in Table~\ref{tab:params}.

\begin{table}[b]
\caption{The SFHo parameter set~\cite{Steiner2013}, as carried by
\texttt{get\_sfho\_nucleonic()}. The CompOSE table states the couplings as
$c_M = g_M/m_M$ in fm and the self-interactions in the dimensionless
$(b, c, \zeta, \xi)$ of Steiner et al.; the values below are those converted
to the $(g, g_2, g_3, c_3, c_4)$ of Eqs.~\eqref{eq:lagrangian}
and~\eqref{eq:Usigma}, with $\sigma$ and $\omega$ measured in MeV. The $\phi$
mass ($1020$~MeV) enters only with hyperons.}
\label{tab:params}
\begin{ruledtabular}
\begin{tabular}{lcc|lcc}
$m_n$ & $939.565346$ & MeV & $g_{\sigma N}$ & $7.53128$ & \\
$m_p$ & $938.272013$ & MeV & $g_{\omega N}$ & $9.02239$ & \\
$m_\sigma$ & $467.4583$ & MeV & $g_{\rho N}$ & $9.30415$ & \\
$m_\omega$ & $782.5011$ & MeV & $g_2$ & $2.95146 \times 10^{3}$ & MeV \\
$m_\rho$ & $763.0001$ & MeV & $g_3$ & $-1.22905 \times 10^{1}$ & \\
 & & & $c_3$ & $-1.78429$ & \\
 & & & $c_4$ & $5.15657$ & \\
\hline
$a_1$ & $-3.81010\times 10^{1}$ & MeV & $b_1$ & $5.51185$ & \\
$a_2$ & $\phantom{-}5.61503\times 10^{-1}$ & & $b_2$ & $-4.62463\times 10^{-5}$ & MeV$^{-2}$ \\
$a_3$ & $\phantom{-}1.45031\times 10^{-3}$ & MeV$^{-1}$ & $b_3$ & $2.81043\times 10^{-7}$ & MeV$^{-4}$ \\
$a_4$ & $\phantom{-}7.11844\times 10^{-5}$ & MeV$^{-2}$ & & & \\
$a_5$ & $\phantom{-}1.60179\times 10^{-7}$ & MeV$^{-3}$ & & & \\
$a_6$ & $\phantom{-}4.05497\times 10^{-10}$ & MeV$^{-4}$ & & & \\
\end{tabular}
\end{ruledtabular}
\end{table}

Non-nucleonic baryons carry their own constant couplings, stored as ratios of
the nucleon ones. The vector sector is SU(6) with an optional overall
enhancement $y$ (the SFHoY choice of Ref.~\cite{Fortin2017}, $y = 1.5$ for
$\Lambda, \Sigma$ and $1.875$ for $\Xi$):
$g_{\omega Y}/g_{\omega N} = y \times \{2/3,\, 2/3,\, 1/3\}$ and
$g_{\phi Y}/g_{\omega N} = y \times \{-\sqrt2/3,\, -\sqrt2/3,\, -2\sqrt2/3\}$
for $\{\Lambda, \Sigma, \Xi\}$, with
$g_{\rho Y}/g_{\rho N} = \{0,\, 2,\, 1\}$. The scalar ratios are either the
published SFHoY values or inverted from the hyperon single-particle potentials
in symmetric matter at saturation,
%
\begin{equation}
U_Y(\nsat) = -g_{\sigma Y}\, \bar\sigma + g_{\omega Y}\, \bar\omega
\qquad (\text{defaults } U_\Lambda = -30,\; U_\Sigma = +30,\;
        U_\Xi = -14 \text{ MeV}),
\label{eq:UY}
\end{equation}
%
one linear equation per hyperon --- and, unlike DD2, with no $\Sigma^R$ in it.
The $\Delta$ quartet has no canonical SFHo coupling table: the default is
universal vector coupling ($g_{\Delta i} = g_{Ni}$), or $g_{\sigma\Delta}$ is
fixed by the same inversion from a chosen $U_\Delta$.

\section{Mean fields and thermodynamics}
\label{sec:fields}

In static uniform matter the meson field equations follow from
Eq.~\eqref{eq:lagrangian}. They are no longer algebraic in the sources, as
they are in DD2: three of the four carry nonlinear terms in the fields
themselves, so the four fields are solved for together with the potentials
rather than being eliminated,
%
\begin{align}
m_\sigma^2 \sigma + g_2 \sigma^2 + g_3 \sigma^3
   - \frac{\partial A}{\partial \sigma}\, \rho^2
   &= (\hbar c)^3 \sum_i g_{\sigma i}\, n^s_i ,
\label{eq:fieldeqs}\\
m_\omega^2 \omega + c_3 \omega^3
   + \frac{\partial A}{\partial \omega}\, \rho^2
   &= (\hbar c)^3 \sum_i g_{\omega i}\, n_i ,
\nonumber\\
m_\rho^2 \rho + c_4 \rho^3 + 2 A\, \rho
   &= (\hbar c)^3 \sum_i g_{\rho i}\, I_{3i}\, n_i ,
\nonumber\\
m_\phi^2 \phi &= (\hbar c)^3 \sum_i g_{\phi i}\, n_i ,
\nonumber
\end{align}
%
with $n_i$ and $n^s_i$ the number and scalar densities of an ideal Fermi gas
of effective mass and effective chemical potential
%
\begin{equation}
\meff{i} = m_i - g_{\sigma i}\, \sigma, \qquad
\mueff{i} = \mu_i - g_{\omega i}\, \omega
          - g_{\rho i}\, I_{3i}\, \rho
          - g_{\phi i}\, \phi .
\label{eq:mueff}
\end{equation}
%
There is no rearrangement term: the couplings do not depend on the state, so
$\Sigma^R = 0$ identically. That is the single structural simplification
relative to DD2, and it is stated in the code rather than omitted. In code, $\mueff{i}$ and $\meff{i}$ are named \texttt{mu\_eff\_i} and
\texttt{m\_eff\_i}.

\subsection{One species as an ideal gas}
\label{sec:kinetic}

Every baryon enters as an ideal Fermi gas of mass $\meff{i}$ at potential
$\mueff{i}$, with degeneracy $g_i$ ($2$ for the octet, $4$ for the $\Delta$)
and antiparticles included. With $E = \sqrt{k^2 + \meff{i}{}^2}$ and
$f_\pm = \big[ 1 + e^{(E \mp \mueff{i})/T} \big]^{-1}$,
%
\begin{align}
n_i &= \frac{g_i}{2\pi^2 (\hbar c)^3} \int_0^\infty \!\! dk\; k^2
       \big( f_+ - f_- \big) ,
\label{eq:nkin}\\
\varepsilon_i &= \frac{g_i}{2\pi^2 (\hbar c)^3} \int_0^\infty \!\! dk\; k^2 E
       \big( f_+ + f_- \big) ,
\label{eq:ekin}\\
P_i &= \frac{g_i}{6\pi^2 (\hbar c)^3} \int_0^\infty \!\! dk\;
       \frac{k^4}{E} \big( f_+ + f_- \big) .
\label{eq:pkin}
\end{align}
%
The remaining two are \emph{not} integrated. The scalar density follows from
the trace of the energy--momentum tensor and the entropy density from the
Euler relation for one species,
%
\begin{equation}
n^s_i = \frac{\varepsilon_i - 3 P_i}{\meff{i}} ,
\qquad
s_i = \frac{\varepsilon_i + P_i - \mueff{i}\, n_i}{T} ,
\label{eq:nsandS}
\end{equation}
%
which is how the code obtains them, and is worth stating because it means an
error in $\varepsilon_i$ or $P_i$ propagates into $n^s_i$ --- and $n^s_i$
sources the $\sigma$ field, so such an error is not confined to the totals.

At $T = 0$ the integrals are elementary. With
$k_{F} = \sqrt{\mueff{i}{}^2 - \meff{i}{}^2}$ (and everything vanishing when
$|\mueff{i}| \le \meff{i}$), writing $L = \ln\big[(k_F + |\mueff{i}|)/\meff{i}\big]$,
%
\begin{align}
n_i &= \mathrm{sgn}(\mueff{i})\, \frac{g_i\, k_F^3}{6\pi^2 (\hbar c)^3} ,
\qquad
n^s_i = \frac{g_i\, \meff{i}}{4\pi^2 (\hbar c)^3}
        \Big( k_F |\mueff{i}| - \meff{i}{}^2 L \Big) ,
\nonumber\\
P_i &= \frac{g_i}{48\pi^2 (\hbar c)^3}
       \Big[ \big( 2k_F^3 - 3\meff{i}{}^2 k_F \big) |\mueff{i}|
             + 3 \meff{i}{}^4 L \Big] ,
\label{eq:T0}\\
\varepsilon_i &= \frac{g_i}{16\pi^2 (\hbar c)^3}
       \Big[ \big( 2k_F^3 + \meff{i}{}^2 k_F \big) |\mueff{i}|
             - \meff{i}{}^4 L \Big] ,
\qquad s_i = 0 .
\nonumber
\end{align}
%
At $T>0$ Eqs.~\eqref{eq:nkin}--\eqref{eq:pkin} are evaluated with the
Johns--Ellis--Lattimer rational approximants~\cite{JohnsEllisLattimer1996},
accurate to $\sim 10^{-4}$, from
\texttt{eos/general/fermi\_integrals} --- the single home of all Fermi and
Bose integrals in this repository, so that no model carries its own. A
Gauss--Laguerre quadrature of the same integrals is available there as the
accuracy reference.

The four source terms of Eq.~\eqref{eq:fieldeqs} are then
%
\begin{equation}
S_\sigma = \sum_i g_{\sigma i}\, n^s_i, \quad
S_\omega = \sum_i g_{\omega i}\, n_i, \quad
S_\rho = \sum_i g_{\rho i}\, I_{3i}\, n_i, \quad
S_\phi = \sum_i g_{\phi i}\, n_i ,
\label{eq:sources}
\end{equation}
%
and the conserved-charge densities are
$\nB = \sum_i B_i n_i$, $n_C = \sum_i C_i n_i + n_C^{\mathrm{mes}}$,
$n_S = \sum_i S_i n_i + n_S^{\mathrm{mes}}$, the meson terms being those of
Sec.~\ref{sec:gas}.

\subsection{The totals}
\label{sec:totals}

The matter block is the baryons plus the mean fields plus any thermal gas,
%
\begin{equation}
\varepsilon = \sum_i \varepsilon_i + \varepsilon_{\mathrm{mf}}
              + \varepsilon_{\mathrm{mes}},
\qquad
P = \sum_i P_i + P_{\mathrm{mf}} + P_{\mathrm{mes}},
\qquad
s = \sum_i s_i + s_{\mathrm{mes}} ,
\label{eq:matter}
\end{equation}
%
with $P_{\mathrm{mf}}$ and $\varepsilon_{\mathrm{mf}}$ from
Eqs.~\eqref{eq:Pmf}--\eqref{eq:epsmf}. Leptons and photons are added on top by
the solver rather than by the matter block, because they are shared by the
whole system rather than belonging to a phase: electrons (and positrons) as an
ideal Fermi gas through the same Eqs.~\eqref{eq:nkin}--\eqref{eq:T0} at
$m = m_e$, $g = 2$; trapped electron neutrinos likewise at $m = 0$, $g = 1$;
and photons as
%
\begin{equation}
P_\gamma = \frac{\pi^2 T^4}{45 (\hbar c)^3}, \qquad
\varepsilon_\gamma = 3 P_\gamma, \qquad
s_\gamma = \frac{4\pi^2 T^3}{45 (\hbar c)^3} .
\label{eq:photons}
\end{equation}
%
The Euler sum reported with the state is
%
\begin{equation}
\sum_i \mu_i n_i
  = \sum_i \big( B_i \muB + C_i \muC + S_i \muS \big) n_i
  + \sum_j \mu^{*}_j n_j
  + \mu_e n_e + \munue n_{\nu_e} ,
\label{eq:eulersum}
\end{equation}
%
with the baryons at their FULL potentials and the meson gas at its EFFECTIVE
ones --- see Sec.~\ref{sec:gas} for why the two differ.

The meson fields contribute to the pressure and the energy density
%
\begin{align}
P_{\mathrm{mf}} &= -V(\sigma)
   + \tfrac12 m_\omega^2 \omega^2 + \tfrac{c_3}{4}\, \omega^4
   + \tfrac12 m_\rho^2 \rho^2 + \tfrac{c_4}{4}\, \rho^4 + A\, \rho^2
   + \tfrac12 m_\phi^2 \phi^2 ,
\label{eq:Pmf}\\
\varepsilon_{\mathrm{mf}} &= +V(\sigma)
   + \tfrac12 m_\omega^2 \omega^2 + \tfrac{3 c_3}{4}\, \omega^4
   + \omega \frac{\partial A}{\partial \omega} \rho^2
   + \tfrac12 m_\rho^2 \rho^2 + \tfrac{3 c_4}{4}\, \rho^4 + A\, \rho^2
   + \tfrac12 m_\phi^2 \phi^2 ,
\label{eq:epsmf}
\end{align}
%
with $V(\sigma) = \tfrac12 m_\sigma^2 \sigma^2 + U(\sigma)$. The asymmetric
terms deserve a word, because two of them were wrong in this code once and
nothing but the Euler identity found it. Eliminating the $\omega$ source
through its own field equation is what turns $\tfrac12 m_\omega^2 \omega^2$
into $\tfrac12 m_\omega^2 \omega^2 + \tfrac{3c_3}{4}\omega^4$; the same
elimination leaves behind $\omega\,(\partial A/\partial\omega)\,\rho^2$,
because $A$ depends on $\omega$ as well as on $\sigma$. Keeping only $b_1$
gives $A = g_\rho^2 b_1 \omega^2$, so that term is $2A\rho^2$ and the cross
coupling enters $\varepsilon$ as $3A\rho^2$ against $A\rho^2$ in $P$ --- the
familiar factor of three of the Horowitz--Piekarewicz
coupling~\cite{HorowitzPiekarewicz2001} this generalises. The $\sigma$
dependence of $A$ has no such partner: $\sigma$ reaches the baryons through
$\meff{}$ and the scalar density rather than through $\mu$, so $V(\sigma)$
simply changes sign between $P$ and $\varepsilon$ and the $a_i$ keep
coefficient one.

\paragraph{The \texttt{thermal\_neutrinos} sector.} With the flag on and
$T > 0$ the solver adds a $\mu = 0$ massless neutrino gas at \textbf{three}
flavours, multiplying the single-flavour result (antiparticles included) into
the totals:
%
\begin{equation}
P \mathrel{+}= 3 P_\nu(0, T) , \qquad
\varepsilon \mathrel{+}= 3 \varepsilon_\nu(0, T) , \qquad
s \mathrel{+}= 3 s_\nu(0, T) .
\label{eq:thermalnu}
\end{equation}
%
It carries no conserved charge and no lepton number, so it touches nothing
else. The factor is exactly $3$, measured: at $\nB = 0.3\ \mathrm{fm}^{-3}$,
$T = 20$~MeV in \texttt{beta\_eq\_neutrinoless}, switching the flag on moves
$P$, $\varepsilon$ and $s$ by $3.000000$ times the single-flavour gas.

Three is right in the modes where the gas is added, and
\texttt{beta\_eq\_neutrino\_trapped} \emph{refuses} the combination rather
than adding it: that mode carries the electron neutrino explicitly at its own
$\munue$, so a three-flavour $\mu = 0$ gas on top would count $\nu_e$ twice,
and the solver raises saying exactly that. In every other mode no neutrino
flavour is tracked in the composition, so all three are untracked.

\textbf{The refusal is a defect, not a design choice, and it is due to be
removed.} The repository conventions state that \texttt{thermal\_neutrinos} is
meaningful alongside \texttt{beta\_eq\_neutrino\_trapped} and that a model
\emph{must not raise} on the combination: the flag is defined by what it does
\emph{not} cover, so under trapping --- where the $e$ and $\mu$ families are
tracked --- it means the $\tau$ family, which is free-streaming and carries no
lepton number the mode constrains. Five models answer this two ways; the three
that succeed are the conformant ones, and SFHo and DID drop the raise. So the
paragraph above states what the code does today and not what it should do:
when the raise goes, the trapped mode gains the two remaining flavours rather
than three. This is not a ledgered gap and there is no ledger entry for it.

Every assembled state is checked against the Hugenholtz--Van Hove
identity~\cite{HugenholtzVanHove1958}
$\varepsilon + P - T s = \sum_i \mu_i n_i$ (sum over every species present,
including leptons and thermal mesons) to a relative $10^{-8}$. That check is
what caught the missing $\omega(\partial A/\partial\omega)\rho^2$ term: it
enters no residual, so the solver converged perfectly on an energy density
wrong by up to $1.8\%$ in asymmetric matter, which put $E_{\mathrm{sym}}$ at
$18.7$~MeV instead of $31.6$ and made $L_{\mathrm{sym}}$ negative.

\section{Conserved charges and species potentials}

The conserved-charge basis is $(B, C, S)$ with $C$ the electric charge of
strongly interacting matter only and $S = +1$ per strange quark ($\Lambda$:
$S=+1$, $\Xi$: $S=+2$; the opposite of the PDG sign, used consistently).
Species potentials are derived, never independent unknowns:
%
\begin{equation}
\mu_i = B_i\, \muB + C_i\, \muC + S_i\, \muS ,
\qquad\text{so}\qquad
\muC = \mu_p - \mu_n,
\quad \text{and beta equilibrium reads } \muC + \mu_e = \munue .
\label{eq:basis}
\end{equation}
%
The muon family is not tracked by this model: $\mu$ appears in no residual and
in no total, and requesting it raises rather than being silently ignored.

\section{The thermal meson gas}
\label{sec:gas}

At $T>0$ the pseudoscalar nonet ($\pi$, $K$, $\eta$, $\eta'$) rides on the
mean field as one-body Bose gases~\cite{Lavagno2010}, with effective potentials tied to
the same vector fields that generate the baryon interactions:
%
\begin{align}
\mu^{*}_{\pi^+} &= \muC - g_{\rho N}\, \rho ,
\nonumber\\
\mu^{*}_{K^+} &= \muC - \muS
   - \big( g_{\omega N} - g_{\omega \Lambda} \big)\, \omega
   - \tfrac12\, g_{\rho N}\, \rho ,
\label{eq:mesonshifts}\\
\mu^{*}_{K^0} &= - \muS
   - \big( g_{\omega N} - g_{\omega \Lambda} \big)\, \omega
   + \tfrac12\, g_{\rho N}\, \rho ,
\nonumber
\end{align}
%
with $\mu^{*}_{j^-} = -\mu^{*}_{j^+}$ and $\mu^{*}_j = 0$ for strangeless
neutral mesons. The $\omega$ shift on the kaons is the light/strange asymmetry
of the vector coupling --- a kaon carries one light quark and one strange
antiquark --- and the $\rho$ shifts follow isospin.

Each species is an ideal Bose gas of its PHYSICAL mass $m_j$: $m_{\pi^\pm} =
139.570$, $m_{\pi^0} = 134.977$, $m_{K^\pm} = 493.677$, $m_{K^0} =
m_{\bar K^0} = 497.611$, $m_\eta = 547.862$, $m_{\eta'} = 957.780$~MeV. The
charged and neutral partners are kept apart rather than averaged: the
splitting is $4.6$~MeV on the pion, and the two carry different charges, so
averaging them would misplace $n_C^{\mathrm{mes}}$ as well as the population.
Nine species in all, since $K^-$ and $\bar K^0$ are tracked separately from
$K^+$ and $K^0$. With $E = \sqrt{k^2 + m_j^2}$, $g_j = 1$, and
$b_\pm = \big[ e^{(E \mp \mu^{*}_j)/T} - 1 \big]^{-1}$,
%
\begin{align}
n_j &= \frac{g_j}{2\pi^2 (\hbar c)^3} \int_0^\infty \!\! dk\; k^2
       \big( b_+ - b_- \big) ,
\qquad
\varepsilon_j = \frac{g_j}{2\pi^2 (\hbar c)^3} \int_0^\infty \!\! dk\; k^2 E
       \big( b_+ + b_- \big) ,
\nonumber\\
P_j &= \frac{g_j}{6\pi^2 (\hbar c)^3} \int_0^\infty \!\! dk\;
       \frac{k^4}{E} \big( b_+ + b_- \big) ,
\qquad
s_j = \frac{\varepsilon_j + P_j - \mu^{*}_j n_j}{T} ,
\label{eq:bose}
\end{align}
%
evaluated with the Bose-case Johns--Ellis--Lattimer
approximants~\cite{JohnsEllisLattimer1996} in
\texttt{eos/general/bose\_integrals}. The gas totals are
$P_{\mathrm{mes}} = \sum_j P_j$, $\varepsilon_{\mathrm{mes}} = \sum_j
\varepsilon_j$, $s_{\mathrm{mes}} = \sum_j s_j$, and its charges
%
\begin{equation}
n_C^{\mathrm{mes}} = \sum_j C_j n_j , \qquad
n_S^{\mathrm{mes}} = \sum_j S_j n_j ,
\label{eq:gascharges}
\end{equation}
%
which are what enter Eq.~\eqref{eq:sources}'s neighbours and the residual rows
of Sec.~\ref{sec:modes}. The gas carries electric
charge and strangeness --- under this sign convention $K^-$ and $\bar K^0$
carry $S=+1$ --- and those densities \emph{enter the equilibrium
constraints}: charge neutrality and the fixed-$Y_C$/$Y_S$ conditions are
imposed on baryons and mesons together. The gas carries no baryon number and
does not source the field equations. Its $P$, $\varepsilon$, $s$ and
$\sum_j \mu^{*}_j n_j$ join the totals and the Hugenholtz--Van Hove sum, the
gas entering at its \emph{effective} potentials because the vector shifts it
carries are sourced by the baryons alone and so have no partner in the field
energy to cancel against.

Bose--Einstein condensation is not implemented, and a state that reaches it is
refused rather than approximated. Every entry point reports
$\max_j |\mu^{*}_j| / m_j$; at $1$ the ideal-gas expressions stop describing
the species, and the solve returns \texttt{converged = False}. In cold
beta-equilibrium nucleonic matter that boundary is crossed just above
saturation, because $\mu^{*}_{\pi^-} \simeq \mu_e$ and $\mu_e$ passes
$m_\pi$ near $\nB \approx 0.27$~fm$^{-3}$: the ideal gas has no repulsive
s-wave $\pi N$ term, which is exactly what suppresses s-wave pion condensation
in real matter. The domain the gas is valid in is the one it was written for,
high $T$ and low $\muB$; negatively charged baryons ($\Sigma^-$, $\Delta^-$)
take over the neutralisation, collapse $\mu_e$, and push the boundary away.

\section{Equilibrium modes and their closures}
\label{sec:modes}

All modes are solved by one residual system over the unknown vector
%
\begin{equation}
x = \big[ \sigma,\ \omega,\ \rho,\ \phi,\ \muB,\ \muC,\
          (\muS),\ (\munue),\ (T) \big] ,
\label{eq:unknowns}
\end{equation}
%
where $\muS$ is present iff strangeness is fixed, $\munue$ iff neutrinos are
trapped, and $T$ iff an entropy per baryon is imposed in its place. All four
fields are always unknowns --- a mean field does not know what is being held
fixed --- and $\muC$ is an unknown in every mode; what the mode changes is the
row that closes it.

The residual is assembled in this order, and these are all of its rows:
%
\begin{align}
R_1 &= \Big[ m_\sigma^2 \sigma + g_2 \sigma^2 + g_3 \sigma^3
        - \tfrac{\partial A}{\partial \sigma}\rho^2
        - (\hbar c)^3 S_\sigma \Big] \big/ \big( m_\sigma^2 \sigma_{\mathrm{sc}} \big) ,
\nonumber\\
R_2 &= \Big[ m_\omega^2 \omega + c_3 \omega^3
        + \tfrac{\partial A}{\partial \omega}\rho^2
        - (\hbar c)^3 S_\omega \Big] \big/ \big( m_\omega^2 \sigma_{\mathrm{sc}} \big) ,
\nonumber\\
R_3 &= \Big[ m_\rho^2 \rho + c_4 \rho^3 + 2 A \rho
        - (\hbar c)^3 S_\rho \Big] \big/ \big( m_\rho^2 \sigma_{\mathrm{sc}} \big) ,
\nonumber\\
R_4 &= \Big[ m_\phi^2 \phi - (\hbar c)^3 S_\phi \Big]
        \big/ \big( m_\phi^2 \sigma_{\mathrm{sc}} \big) ,
\label{eq:residual}\\
R_5 &= \sum_i B_i n_i - \nB^{\mathrm{target}} ,
\nonumber\\
R_6 &= \begin{cases}
  \sum_i C_i n_i + n_C^{\mathrm{mes}} - n_e & \text{$C$ equilibrated} \\
  \sum_i C_i n_i + n_C^{\mathrm{mes}} - Y_C\, \nB & \text{$Y_C$ imposed}
\end{cases}
\nonumber\\
R_7 &= \sum_i S_i n_i + n_S^{\mathrm{mes}} - Y_S\, \nB
       \qquad \text{(iff $Y_S$ imposed)},
\nonumber\\
R_8 &= \big( n_e + n_{\nu_e} \big) / \nB - Y_{L_e}
       \qquad \text{(iff the electron family is trapped)},
\nonumber\\
R_9 &= s_{\mathrm{tot}} / \nB - (S/A)
       \qquad \text{(iff an entropy per baryon replaces $T$)}.
\nonumber
\end{align}
%
The sources $S_M$ are Eq.~\eqref{eq:sources}, the gas charges
Eq.~\eqref{eq:gascharges}, and $\sigma_{\mathrm{sc}} = 30$~MeV is a typical scalar field
at saturation, so that the four field rows are dimensionless and of order
unity rather than of order $m_M^2 \sigma \sim 10^7$~MeV$^3$ --- without that
scaling they would dominate the residual norm and the charge rows would never
be converged. Rows $R_1$--$R_6$ are always present; the last three appear with
their conditions. Per mode:
%
\begin{center}
\begin{tabular}{llll}
\toprule
mode & independent variables & closing constraints & potentials \\
\midrule
\texttt{beta\_eq\_neutrinoless} & $(\nB, T)$ &
  $\sum_i C_i n_i + n_C^{\mathrm{mes}} - n_e = 0$ &
  $\muS = \munue = 0$ \\
\texttt{beta\_eq\_neutrino\_trapped} & $(\nB, Y_{L_e}, T)$ &
  neutrality; $(n_e + n_{\nu_e})/\nB = Y_{L_e}$ &
  $\muS = 0$; $\munue$ unknown \\
\texttt{fixed\_YC} & $(\nB, Y_C, T)$ &
  $\big(\textstyle\sum_i C_i n_i + n_C^{\mathrm{mes}}\big)/\nB = Y_C$ &
  $\muS = \munue = 0$ \\
\texttt{fixed\_YC\_YS} & $(\nB, Y_C, Y_S, T)$ &
  the $Y_C$ row; $\big(\textstyle\sum_i S_i n_i + n_S^{\mathrm{mes}}\big)/\nB = Y_S$ &
  $\muS$ unknown; $\munue = 0$ \\
\bottomrule
\end{tabular}
\end{center}
%
In the fixed-fraction modes the flag \texttt{leptons} selects between the
strongly-interacting slice alone (electrically charged; what a mixed-phase
construction consumes) and the same slice plus neutralizing electrons
($n_e = Y_C \nB$) --- the leptons do not source the fields and enter no
residual row, so they close post hoc through a single 1-D root. Wherever a
temperature is accepted, entropy per baryon $s/\nB$ is accepted in its place;
here $T$ joins the unknown vector and $s/\nB - (S/A) = 0$ joins the rows,
rather than being an outer 1-D solve.

Non-convergence is a return value at every layer, judged on the largest
residual and reported on the solved record, never an exception --- a sampler
must be able to score a bad point and move on.

\section{Nuclear-matter parameters}

The forward map evaluates, at the parametrization's own saturation density
(the $P=0$ root of symmetric matter): $E_{\mathrm{sat}} = E/A$, $m^*/m$,
$K_{\mathrm{sat}} = 9 \nsat^2\, \partial^2_n (E/A)$,
$Q_{\mathrm{sat}} = 27 \nsat^3\, \partial^3_n (E/A)$, and in the isovector
sector the mean-field closed form of Steiner et al.~\cite{Steiner2005}
Eq.~(20),
%
\begin{equation}
E_{\mathrm{sym}}(\nB) = \frac{k_F^2}{6 E^*_F}
  + \frac{\nB\, g_{\rho N}^2}{8 \big( m_\rho^2 + 2 A \big)} ,
\qquad L_{\mathrm{sym}} = 3 \nsat\, E_{\mathrm{sym}}'(\nsat),
\qquad K_{\mathrm{sym}} = 9 \nsat^2\, E_{\mathrm{sym}}''(\nsat) ,
\label{eq:esym}
\end{equation}
%
where the $2A$ in the denominator is the whole point of the cross coupling: it
is a function of the isoscalar fields and therefore of density, so it bends
$E_{\mathrm{sym}}(\nB)$ without moving its value at saturation. The map
returns the same key set as \texttt{eos.dd2.compute\_nmp}, so one caller reads
either model, and it reproduces the published SFHo values --- $\nsat = 0.158$,
$E_{\mathrm{sat}} = -16.2$, $E_{\mathrm{sym}} = 31.6$, $L_{\mathrm{sym}} =
47.1$~MeV --- to better than $0.3\%$. The verification suite computes
$E_{\mathrm{sym}}$ a second, independent way, as the $\delta^2$ curvature of
$E/A$ in the isospin asymmetry, and requires the two to agree: the closed form
is a $\rho$-field response and the curvature is a property of $\varepsilon$,
so an error in the energy density moves one and not the other.

The inverse map is triangular, and that is a property of the model rather than
a solver tactic: in symmetric matter $\rho$ and $A(\sigma,\omega)\rho^2$ drop
out of every equation of Sec.~\ref{sec:fields}, so the isoscalar sector does
not see the isovector couplings at all and is solved first.

Its isoscalar half is the classical Boguta--Bodmer
inversion~\cite{BogutaBodmer1977}: four conditions $\{P(\nsat) = 0,
E_{\mathrm{sat}}, m^*/m, K_{\mathrm{sat}}\}$ against four unknowns
$\{g_{\sigma N}, g_{\omega N}, g_2, g_3\}$ at fixed $m_\sigma$, $m_\omega$,
$c_3$, with no structural closure needed. It is solved in the reduced
self-couplings the published table states, $g_2 = b\, m_N g_{\sigma N}^3$ and
$g_3 = c\, g_{\sigma N}^4$, because $b \sim 7\times10^{-3}$ and
$c \sim -4\times10^{-3}$ sit beside couplings of order ten where $g_2$ and
$g_3$ span $3\times10^3$~MeV and $-12$.

Its isovector half needs a choice: two conditions $\{E_{\mathrm{sym}},
L_{\mathrm{sym}}\}$ face $g_{\rho N}$ plus the nine shape coefficients of
Eq.~\eqref{eq:Afunction}, so exactly two of the ten are freed and the rest
pinned. That choice is physics rather than bookkeeping, because it decides how
$E_{\mathrm{sym}}$ behaves above saturation where no nuclear-matter parameter
constrains it. The closure used is $(g_{\rho N}, b_1)$, for three reasons. It
is the best conditioned of the candidates --- the $2\times2$ Jacobian in
logarithmic knobs at the SFHo point has condition number $3.40$, against
$3.53$ for an overall scale on $f$ and $11.60$ for $a_1$, which moves
$L_{\mathrm{sym}}$ by less than $2$~MeV per $e$-fold. It has by far the widest
reach: scanning both knobs at $E_{\mathrm{sym}}$ held near $31.5$~MeV, the
accessible $L_{\mathrm{sym}}$ spans $[-6, 146]$~MeV, against $[-34, 59]$ for
the scale and $[27, 69]$ for $a_1$. And $b_1$ IS the Horowitz--Piekarewicz
$\Lambda_v \omega^2\rho^2$ coupling~\cite{HorowitzPiekarewicz2001}, the
standard way this family tunes $L_{\mathrm{sym}}$ at fixed
$E_{\mathrm{sym}}$, so an inverted parameter set stays comparable to published
ones. Note that the overall scale, which looks least invasive because it
preserves the \emph{shape} of $A$, is the most invasive to the physics:
scaling $A$ changes how $2A$ competes with $m_\rho^2$ in
Eq.~\eqref{eq:esym} as density rises, and at $L_{\mathrm{sym}} = 70$~MeV it
puts $E_{\mathrm{sym}}(3\nsat)$ at $55.8$~MeV against $52.2$ for $b_1$ and
$49.6$ for $a_1$.

$Q_{\mathrm{sat}}$ and $K_{\mathrm{sym}}$ are then \emph{predictions} of the
closure and are reported as such. One caveat belongs in the specification: the
isovector system has a second branch. The potential term of
Eq.~\eqref{eq:esym} saturates at $\nB/(16 f)$ as $g_{\rho N}$ grows, because
$A = g_{\rho N}^2 f$ carries the same coupling in the denominator, so a
runaway $(g_{\rho N}, b_1)$ can reproduce a target set exactly --- at low
$m^*/m$ and low $L_{\mathrm{sym}}$ it is the only root there is. The inversion
therefore checks $|2A| < m_\rho^2$ at saturation after converging, which is
the assumption the model form is written under, and reports a fit that lands
outside it as a failure rather than returning it. Published SFHo sits at
$2A/m_\rho^2 = +0.37$ and every accepted fit from $L_{\mathrm{sym}} = 40$ to
$140$~MeV inside $[-0.40, +0.69]$.

\section{Numerical implementation}

The reference path is plain NumPy/SciPy: the residual above, with MINPACK
building its own forward-difference Jacobian. It is the correctness oracle and
is never bypassed. Alongside it, \texttt{backends/} carries a hand-derived
analytic Jacobian of the same residual, agreeing with a central difference of
it to better than $10^{-7}$ across every mode; its kinetic derivatives
are closed forms at $T=0$ and per-species central differences of the JEL
integrals at $T>0$, since those expose no exact derivative. It is not selected
by default, because it is not faster here --- it cuts the residual evaluations
of a warm-started sweep by about a third but costs more per evaluation than it
saves. What it is for is the susceptibility matrix
$\chi_{ab} = \partial n_a / \partial \mu_b$ for $a,b \in (B, C, S)$, which is
not a finite-difference quantity at all: the solver never varies the three
potentials independently, so there is no sequence to difference along. It is
read off the Jacobian through the field-response identity
%
\begin{equation}
\frac{\partial n_a}{\partial \mu_b}
  = \left.\frac{\partial n_a}{\partial \mu_b}\right|_{F}
  + \frac{\partial n_a}{\partial F}\,
    \left( -\Big(\frac{\partial G}{\partial F}\Big)^{-1}
             \frac{\partial G}{\partial \mu_b} \right),
\label{eq:chi}
\end{equation}
%
with $F$ the four mean fields and $G$ their field equations, and it is checked
against the inverse map $\partial\mu_a/\partial n_b$, which the
\texttt{fixed\_YC\_YS} mode supplies without touching the Jacobian at all.

Density sweeps are warm-started: each solved point seeds the next. Two
degeneracies are known and are properties of the model rather than of the
solver. Where a conserved charge is carried by no populated species its
potential is unpinned --- the strangeness row is flat in $\muS$ when
$n_S \to 0$, and $\mu_e$ is only weakly determined at $Y_C = 0$ --- so
neither quantity is frozen in the regression baseline where its conjugate
density vanishes, and the susceptibility check is placed where strangeness is
populated.

\section{The published parameter sets}
\label{sec:sets}

\texttt{Parameters.default()} returns \texttt{SFHo\_Nucleonic};
\texttt{Parameters.named(key)} returns any of the five, and
\texttt{PUBLISHED\_SETS} --- exported from the package and from
\texttt{eos.sfho.parameters} --- maps the keys to builders. The keys are
public API:
%
\begin{center}
\begin{tabular}{ll}
\toprule
key & contents \\
\midrule
\texttt{SFHo\_Nucleonic} & nucleons only, the CompOSE SFHo table. What
   \texttt{default()} returns, what the \\
 & nuclear-matter map reports the published values against, and what \\
 & the regression baseline is frozen at \\
\texttt{SFHoY\_Fortin} & $+$ hyperons, scaled vector couplings ($y = 1.5$),
   scalar couplings \\
 & from $U_Y$~\cite{Fortin2017} \\
\texttt{SFHoY*\_Fortin} & $+$ hyperons, SU(6) vector couplings (same
   reference) \\
\texttt{SFHo\_2fam\_phi} & $+$ hyperons and $\Delta$s, SU(6) vectors,
   hyperons coupled to $\phi$ \\
\texttt{SFHo\_2fam} & as \texttt{SFHo\_2fam\_phi} with $g_\phi = 0$ for every
   strange baryon \\
\bottomrule
\end{tabular}
\end{center}
%
All five share the isoscalar and isovector couplings of
Sec.~\ref{sec:lagrangian}; they differ only in which baryons are active and
how those baryons couple.

\section{The phase-adapter surface}
\label{sec:adapter}

What \texttt{eos/mixed} consumes across the phase-adapter contract, and the
layer beneath it:
%
\begin{quote}
\texttt{thermo\_from\_mu(par, flags, mu\_B, mu\_C, mu\_S=0.0, T=0.0,\\
\phantom{thermo\_from\_mu(}n\_B\_guess=0.2, x0=None, x0\_fallback=None,\\
\phantom{thermo\_from\_mu(}return\_state=False)}
\end{quote}
%
returns a \texttt{PhaseThermo}, or the pair
$(\texttt{PhaseThermo}, \{\texttt{"x\_phase"}: [\sigma, \omega, \rho, \phi]\})$
with \texttt{return\_state}. It solves the four meson fields against the field
residual at the \emph{given} potentials and then assembles, and raises
\texttt{RuntimeError} when no guess meets the per-row scaled residual
tolerance --- an internal layer may raise; the public entry points catch and
report.

It takes the \textbf{physical} $\muB$. SFHo's couplings are constants, so
there is no rearrangement term and no kinetic potential --- unlike the DD2
counterpart, which takes the kinetic one. Which of the two a phase's slot
carries is a declared property of the phase, never an engine assumption. That
counterpart is also still spelled \texttt{thermo\_at\_potentials}: seven of
the ten models use \texttt{thermo\_from\_mu} for this surface, SFHo is now one
of them, and DD2 is the outstanding one.
%
\begin{quote}
\texttt{thermo\_from\_fields(mu\_B, mu\_C, mu\_S, sigma, omega, rho, phi, T,\\
\phantom{thermo\_from\_fields(}particles, params,\\
\phantom{thermo\_from\_fields(}include\_pseudoscalar\_mesons=False)}
\end{quote}
%
is the lower layer, which additionally takes the solved mean fields --- the
name says what the function takes, and that is the whole distinction between
the two. Matter only: baryons, the mean fields, and any thermal $\pi/K/\eta$
gas. No leptons and no photons; the solver adds those, because they are shared
by the whole system rather than belonging to a phase.

\section{What \texttt{eos\_response} returns}
\label{sec:response}

\texttt{frozen='equilibrium'} is the implemented freeze. It returns
%
\begin{center}
\begin{tabular}{lll}
\toprule
key & quantity & availability \\
\midrule
\texttt{cs2\_isothermal} & $(\partial P/\partial\varepsilon)_T$ along the
                           mode's sequence & always \\
\texttt{chi}             & $\chi_{ab} = \partial n_a/\partial\mu_b$,
                           Eq.~\eqref{eq:chi}, from the Jacobian & always \\
\texttt{cs2\_adiabatic}  & the same derivative at fixed $s/\nB$ & $T > 0$
                           only \\
\texttt{C\_V}            & $(T/\nB)(\partial s/\partial T)_{\nB}$ & $T > 0$
                           only \\
\texttt{C\_P}            & $(T/\nB)(\partial s/\partial T)_{P}$ & $T > 0$
                           only \\
\texttt{Gamma\_th}       & the thermal index & $T > 0$ only \\
\bottomrule
\end{tabular}
\end{center}
%
all but $\chi$ by finite differences along re-solved sequences. \textbf{The two
sound speeds are named for the thermal condition they are taken at}, because
at $T > 0$ they are different numbers --- they differ by $C_P/C_V$ --- and a
bare \texttt{cs2} whose meaning depended on the arguments would be the more
misleading of the two. Every other freeze of the specification --- frozen
per-species composition, frozen conserved fractions, the leptonic
re-neutralization variants --- raises \texttt{NotImplementedError} naming
itself.

\bibliography{../../docs/eos}

\end{document}
